\textbf{Phần 1. Trắc nghiệm lựa chọn}
\setcounterpageref{ex}{0}

\begin{ex} %[3D3V31]
Điền số thích hợp vào chỗ chấm: $1\text{ kg} = \dots\text{ g}$.
\choice
{$100\text{ g}$}
{$10\text{ g}$}
{\True $1000\text{ g}$}
{$10000\text{ g}$}
\loigiai{
$1\text{ kg} = 1000\text{ g}$.
\textbf{Đáp án của câu hỏi là C}
}
\end{ex}

\begin{ex} %[3D3V31]
Một quả táo nặng $180\text{ g}$, quả lê nặng hơn quả táo $45\text{ g}$. Cả hai quả nặng tất cả là
\choice
{\True $405\text{ g}$}
{$360\text{ g}$}
{$420\text{ g}$}
{$395\text{ g}$}
\loigiai{
Quả lê nặng: $180 + 45 = 225\text{ g}$. Cả hai quả nặng: $180 + 225 = 405\text{ g}$.
\textbf{Đáp án của câu hỏi là A}
}
\end{ex}

\begin{ex} %[3D3V31]
Trên đĩa cân đĩa thứ nhất đặt một gói đường và một quả cân $200\text{ g}$. Đĩa cân thứ hai đặt một quả cân $1\text{ kg}$. Biết cân thăng bằng, gói đường nặng là
\choice
{$1200\text{ g}$}
{$700\text{ g}$}
{$900\text{ g}$}
{\True $800\text{ g}$}
\loigiai{
Đổi $1\text{ kg} = 1000\text{ g}$. Gói đường nặng: $1000 - 200 = 800\text{ (g)}$.
\textbf{Đáp án của câu hỏi là D}
}
\end{ex}

\begin{ex} %[3D3V31]
Một hộp bánh quy gồm vỏ hộp nặng $80\text{ g}$ và 4 gói bánh nhỏ bên trong, mỗi gói bánh nặng $120\text{ g}$. Cả hộp bánh cân nặng là
\choice
{$480\text{ g}$}
{\True $560\text{ g}$}
{$600\text{ g}$}
{$520\text{ g}$}
\loigiai{
Bốn gói bánh nặng: $120 \times 4 = 480\text{ g}$. Cả hộp nặng: $480 + 80 = 560\text{ (g)}$.
\textbf{Đáp án của câu hỏi là B}
}
\end{ex}

\begin{ex} %[3D3V31]
Mẹ có $1\text{ kg}$ bột mì. Mẹ dùng $350\text{ g}$ để làm bánh rán và $450\text{ g}$ để làm bánh bao. Mẹ còn lại bao nhiêu gam bột mì?
\choice
{$150\text{ g}$}
{$250\text{ g}$}
{\True $200\text{ g}$}
{$300\text{ g}$}
\loigiai{
Đổi $1\text{ kg} = 1000\text{ g}$. Số bột đã dùng: $350 + 450 = 800\text{ g}$. Còn lại: $1000 - 800 = 200\text{ (g)}$.
\textbf{Đáp án của câu hỏi là C}
}
\end{ex}

\begin{ex} %[3D3V31]
Có 5 túi muối giống nhau cân nặng tất cả $750\text{ g}$. Hỏi 8 túi muối như thế cân nặng bao nhiêu gam?
\choice
{\True $1200\text{ g}$}
{$1000\text{ g}$}
{$1500\text{ g}$}
{$1150\text{ g}$}
\loigiai{
Mỗi túi muối nặng: $750 : 5 = 150\text{ g}$. Tám túi nặng: $150 \times 8 = 1200\text{ (g)}$.
\textbf{Đáp án của câu hỏi là A}
}
\end{ex}

\begin{ex} %[3D3V31]
Một quả dưa hấu nặng $2\text{ kg } 500\text{ g}$. Người ta bổ quả dưa thành 5 miếng đều nhau. Mỗi miếng dưa hấu nặng là
\choice
{$400\text{ g}$}
{$600\text{ g}$}
{$450\text{ g}$}
{\True $500\text{ g}$}
\loigiai{
Đổi $2\text{ kg } 500\text{ g} = 2500\text{ g}$. Mỗi miếng dưa nặng: $2500 : 5 = 500\text{ (g)}$.
\textbf{Đáp án của câu hỏi là D}
}
\end{ex}

\begin{ex} %[3D3V31]
Một gói mì tôm nặng $75\text{ g}$. Hỏi một thùng gồm 30 gói mì tôm như thế cân nặng bao nhiêu ki-lô-gam và gam (bỏ qua khối lượng vỏ thùng)?
\choice
{$225\text{ g}$}
{\True $2\text{ kg } 250\text{ g}$}
{$2\text{ kg } 500\text{ g}$}
{$3\text{ kg}$}
\loigiai{
Tổng cân nặng: $75 \times 30 = 2250\text{ g} = 2\text{ kg } 250\text{ g}$.
\textbf{Đáp án của câu hỏi là B}
}
\end{ex}

\begin{ex} %[3D3V31]
Cô Lan mua $1\text{ kg}$ thịt bò. Cô chia đều số thịt đó vào 4 hộp để trong tủ lạnh. Mỗi hộp thịt bò nặng là
\choice
{$200\text{ g}$}
{$300\text{ g}$}
{\True $250\text{ g}$}
{$400\text{ g}$}
\loigiai{
Đổi $1\text{ kg} = 1000\text{ g}$. Mỗi hộp nặng: $1000 : 4 = 250\text{ (g)}$.
\textbf{Đáp án của câu hỏi là C}
}
\end{ex}

\begin{ex} %[3D3V31]
Tìm số quả cân $50\text{ g}$ cần dùng để thăng bằng với một quả cân $200\text{ g}$ và một quả cân $100\text{ g}$.
\choice
{\True $6$ quả cân}
{$5$ quả cân}
{$4$ quả cân}
{$8$ quả cân}
\loigiai{
Tổng khối lượng cần cân là: $200 + 100 = 300\text{ g}$. Số quả cân $50\text{ g}$ cần là: $300 : 50 = 6$ quả.
\textbf{Đáp án của câu hỏi là A}
}
\end{ex}

\begin{ex} %[3D3V31]
Tính giá trị của biểu thức: $450\text{ g} : 9 + 85\text{ g} \times 2$.
\choice
{$200\text{ g}$}
{$210\text{ g}$}
{$230\text{ g}$}
{\True $220\text{ g}$}
\loigiai{
$450 : 9 = 50\text{ g}$; $85 \times 2 = 170\text{ g}$; $50 + 170 = 220\text{ g}$.
\textbf{Đáp án của câu hỏi là D}
}
\end{ex}

\begin{ex} %[3D3V31]
Một con gà mái nặng $1\text{ kg } 800\text{ g}$. Con vịt nhẹ hơn con gà mái $300\text{ g}$. Cả gà mái và vịt cân nặng là
\choice
{$3\text{ kg}$}
{\True $3\text{ kg } 300\text{ g}$}
{$3\text{ kg } 500\text{ g}$}
{$2\text{ kg } 800\text{ g}$}
\loigiai{
Con vịt nặng: $1800 - 300 = 1500\text{ g}$. Cả hai con nặng: $1800 + 1500 = 3300\text{ g} = 3\text{ kg } 300\text{ g}$.
\textbf{Đáp án của câu hỏi là B}
}
\end{ex}

\textbf{Phần 2. Trắc nghiệm đúng sai}
\setcounterpageref{ex}{0}

\begin{ex} %[3D3V31]
Xét tính đúng sai của các mối liên hệ về đơn vị đo gam:
\choiceTF
{\True Gam được viết tắt là $\text{g}$.}
{$1\text{ kg} = 100\text{ g}$.}
{\True Nửa ki-lô-gam bằng $500\text{ g}$.}
{$450\text{ g} + 650\text{ g} = 1\text{ kg}$.}
\loigiai{
\begin{itemize}
    \item Ý a) Đúng: Kí hiệu chuẩn.
    \item Ý b) Sai: Sai vì $1\text{ kg} = 1000\text{ g}$.
    \item Ý c) Đúng: $1000 : 2 = 500\text{ g}$.
    \item Ý d) Sai: Sai vì tổng là $1100\text{ g} = 1\text{ kg } 100\text{ g}$.
\end{itemize}
\begin{itemize}
    \item a) \textbf{Đúng}.
    \item b) \textbf{Sai}.
    \item c) \textbf{Đúng}.
    \item d) \textbf{Sai}.
\end{itemize}
}
\end{ex}

\begin{ex} %[3D3V31]
Một túi gạo có ghi khối lượng tịnh $1\text{ kg}$:
\choiceTF
{\True Khối lượng gạo thực tế bên trong túi không tính vỏ là $1000\text{ g}$.}
{\True Nếu vỏ bao bì nặng $15\text{ g}$ thì cả túi gạo và vỏ nặng $1015\text{ g}$.}
{Nếu bớt ra khỏi túi $300\text{ g}$ gạo thì trong túi còn lại $800\text{ g}$ gạo.}
{\True Hai túi gạo như thế có tổng khối lượng gạo là $2000\text{ g}$.}
\loigiai{
\begin{itemize}
    \item Ý a) Đúng: Định nghĩa khối lượng tịnh.
    \item Ý b) Đúng: Tổng khối lượng cả bì.
    \item Ý c) Sai: Sai vì còn lại $1000 - 300 = 700\text{ g}$.
    \item Ý d) Đúng: $1000 \times 2 = 2000\text{ g}$.
\end{itemize}
\begin{itemize}
    \item a) \textbf{Đúng}.
    \item b) \textbf{Đúng}.
    \item c) \textbf{Sai}.
    \item d) \textbf{Đúng}.
\end{itemize}
}
\end{ex}

\textbf{Phần 3. Trả lời ngắn}
\setcounterpageref{ex}{0}

\begin{ex} %[3D3V31]
Tính kết quả phép tính: $1\text{ kg} - 350\text{ g}$ theo đơn vị gam.
\shortans{650}
\loigiai{
\textbf{Đáp số:} $650$\\
$1000 - 350 = 650\text{ (g)}$.
}
\end{ex}

\begin{ex} %[3D3V31]
Một gói kẹo nặng $125\text{ g}$. Hỏi 4 gói kẹo như thế nặng bao nhiêu gam?
\shortans{500}
\loigiai{
\textbf{Đáp số:} $500$\\
$125 \times 4 = 500\text{ (g)}$.
}
\end{ex}

\begin{ex} %[3D3V31]
Người ta chia $600\text{ g}$ đường vào 3 lọ bằng nhau. Hỏi mỗi lọ đựng bao nhiêu gam đường?
\shortans{200}
\loigiai{
\textbf{Đáp số:} $200$\\
$600 : 3 = 200\text{ (g)}$.
}
\end{ex}

\begin{ex} %[3D3V31]
Có hai gói chè, gói thứ nhất nặng $150\text{ g}$, gói thứ hai nặng gấp đôi gói thứ nhất. Hỏi cả hai gói chè nặng bao nhiêu gam?
\shortans{450}
\loigiai{
\textbf{Đáp số:} $450$\\
Gói hai: $300\text{ g}$. Cả hai: $150 + 300 = 450\text{ (g)}$.
}
\end{ex}

\begin{ex} %[3D3V31]
Một hộp sữa bột nặng $900\text{ g}$. Sau khi dùng một số gam sữa thì lượng sữa còn lại trong hộp giảm đi 3 lần. Hỏi trong hộp còn lại bao nhiêu gam sữa (bỏ qua vỏ)?
\shortans{300}
\loigiai{
\textbf{Đáp số:} $300$\\
$900 : 3 = 300\text{ (g)}$.
}
\end{ex}

\begin{ex} %[3D3V31]
Cần bao nhiêu quả cân $200\text{ g}$ để có tổng khối lượng đúng bằng $1\text{ kg}$?
\shortans{5}
\loigiai{
\textbf{Đáp số:} $5$\\
$1000 : 200 = 5$ quả cân.
}
\end{ex}
